\textbf{1. Kiến thức cần nhớ}
\begin{itemize}
\item \textbf{Giao của hai tập hợp $A$ và $B$:} Gồm các phần tử thuộc cả $A$ và $B$.
\[ \boxed{A \cap B = \{x \mid x \in A \text{ và } x \in B\}} \]
\item \textbf{Hợp của hai tập hợp $A$ và $B$:} Gồm các phần tử thuộc $A$ hoặc thuộc $B$.
\[ \boxed{A \cup B = \{x \mid x \in A \text{ hoặc } x \in B\}} \]
\item \textbf{Hiệu của hai tập hợp $A$ và $B$:} Gồm các phần tử thuộc $A$ nhưng không thuộc $B$.
\[ \boxed{A \setminus B = \{x \mid x \in A \text{ và } x \notin B\}} \]
\item \textbf{Phần bù:} Nếu $A$ là tập con của $E$ ($A \subset E$) thì hiệu $E \setminus A$ gọi là phần bù của $A$ trong $E$, kí hiệu $\boxed{C_E A = E \setminus A}$.
\item \textbf{Công thức số phần tử} (với $A, B$ là các tập hữu hạn):
\[ \boxed{n(A \cup B) = n(A) + n(B) - n(A \cap B)} \]
Nếu $A \cap B = \emptyset$ thì $n(A \cup B) = n(A) + n(B)$.
\end{itemize}

\textbf{2. Phương pháp giải các dạng}

%% DANG 0D1?3-1
\textbf{Dạng 1. Giao và hợp của hai tập hợp (rời rạc)}
\begin{enumerate}
\item Nếu tập hợp cho bằng tính chất đặc trưng, liệt kê tất cả các phần tử của từng tập hợp $A, B$.
\item Tìm $A \cap B$: chọn các phần tử xuất hiện ở cả $A$ và $B$.
\item Tìm $A \cup B$: gộp tất cả các phần tử của $A$ và $B$ (mỗi phần tử chỉ viết một lần).
\end{enumerate}
Lưu ý: Không nhầm lẫn kí hiệu giao ($\cap$, lấy phần chung) và hợp ($\cup$, lấy tất cả); nếu không có phần tử chung thì $A \cap B = \emptyset$.

%% DANG 0D1?3-2
\textbf{Dạng 2. Hiệu và phần bù của hai tập hợp (rời rạc)}
\begin{enumerate}
\item Liệt kê các phần tử của các tập hợp $A, B$ (hoặc $E, A$).
\item Tìm $A \setminus B$: lấy các phần tử thuộc $A$ và gạch bỏ những phần tử có mặt trong $B$.
\item Tìm $C_E A$: kiểm tra điều kiện $A \subset E$, sau đó tính $E \setminus A$.
\end{enumerate}
Lưu ý: Phép hiệu không có tính giao hoán, nói chung $A \setminus B \neq B \setminus A$; chỉ tồn tại $C_E A$ khi $A \subset E$.

%% DANG 0D1?3-3
\textbf{Dạng 3. Giao và hợp (dùng đoạn, khoảng)}
\begin{enumerate}
\item Biểu diễn các tập hợp $A, B$ trên cùng một trục số (hoặc hai trục số thẳng hàng) theo đúng thứ tự các đầu mút.
\item Tìm $A \cap B$: gạch bỏ phần không thuộc $A$ và phần không thuộc $B$; phần không bị gạch chính là $A \cap B$.
\item Tìm $A \cup B$: tô đậm cả hai tập $A$ và $B$; toàn bộ phần được tô đậm chính là $A \cup B$.
\end{enumerate}
Lưu ý: Cẩn thận dấu ngoặc vuông $[$ $]$ (lấy đầu mút) và ngoặc tròn $($ $)$ (không lấy đầu mút) tại các điểm biên trùng nhau.

%% DANG 0D1?3-4
\textbf{Dạng 4. Hiệu và phần bù (dùng đoạn, khoảng)}
\begin{enumerate}
\item Biểu diễn tập hợp $A$ trên trục số (gạch bỏ phần ngoài $A$).
\item Biểu diễn tập hợp $B$ và gạch bỏ toàn bộ phần thuộc $B$; phần còn lại không bị gạch là $A \setminus B$.
\item Với phần bù $C_{\mathbb{R}} A = \mathbb{R} \setminus A$: bỏ đi phần thuộc $A$ trên trục số thực $\mathbb{R}$.
\end{enumerate}
Lưu ý: Khi gạch bỏ một khoảng mở $(a; b)$, các đầu mút $a, b$ vẫn được giữ lại (đổi thành ngoặc vuông); ngược lại, khi gạch bỏ đoạn $[a; b]$, các đầu mút $a, b$ bị loại đi (đổi thành ngoặc tròn).

%% DANG 0D1?3-5
\textbf{Dạng 5. Toán thực tế ứng dụng của tập hợp}
\begin{enumerate}
\item Gọi tên các tập hợp biểu thị từng nhóm đối tượng trong đề bài và vẽ biểu đồ Ven minh họa.
\item Xác định số phần tử của từng tập hợp, phần giao hoặc phần hợp từ giả thiết.
\item Áp dụng công thức $n(A \cup B) = n(A) + n(B) - n(A \cap B)$ (hoặc đặt ẩn theo các miền rời nhau trên biểu đồ Ven) để tìm đại lượng cần tính.
\end{enumerate}
Lưu ý: Chú ý tách riêng những đối tượng không thuộc bất kì tập hợp nào (nằm ngoài $A \cup B$) trước khi áp dụng công thức.
